\section{Methods}\label{sec:methods}

\subsection{Pipeline overview}

The generator produced synthetic American option prices through a three-stage pipeline:
\begin{align}\label{eq:pipeline}
    &\underbrace{\text{JumpHMM}(\text{prices})}_{\text{Stage 1}}
    \;\to\; \{S_t,\, s_t,\, \text{jumps}\}
    \;\to\; \underbrace{dv \text{ process}}_{\text{Stage 2}}
    \;\to\; \sigma_{\text{imp}} = \sqrt{v_t} \nonumber\\
    &\quad\to\; \underbrace{\text{CRR}(S_t,\, \sigma_{\text{imp}},\, K,\, \text{DTE},\, r_f)}_{\text{Stage 3}}
    \;\to\; P_{\text{american}}
\end{align}
This ordering was dictated by data flow: the JumpHMM produced the price paths and regime-state sequences that the Heston process required as inputs, and the Heston process produced the IV values that the CRR tree required for pricing. Each stage was self-contained, so the pipeline could be run forward without iteration or feedback from downstream stages.

\subsection{Modified Heston variance process}\label{sec:theta_function}

We modified the standard Heston process~\eqref{eq:heston_standard} by making the mean-reversion target time-varying:
\begin{equation}\label{eq:heston_modified}
    dv = \kappa\bigl(\theta(t) - v\bigr)\,dt + \sigma_v \sqrt{v}\,dW_v
\end{equation}
The mean-reversion target decomposed into a regime-dependent base level, a market mood modulator, and a contract-specific smile and term-structure adjustment:
\begin{equation}\label{eq:theta}
    \theta(t) = \theta_{s_t} \cdot \bigl(1 + \gamma \cdot M_t\bigr) \cdot \psi(\mathrm{DTE}_t,\; K/S_t)
\end{equation}
where $\theta_{s_t}$ was a learned variance level for HMM state $s_t$ (one value per state, $s_t \in \{1, \ldots, N\}$); $\gamma \geq 0$ was the market mood sensitivity controlling how strongly aggregate stress elevated IV; $M_t \in [0, 1]$ was the aggregate market mood; and $\psi(\mathrm{DTE}, K/S)$ was the smile and term-structure adjustment. The regime-dependent base $\theta_{s_t}$ linked the Heston process directly to the JumpHMM state sequence: when the HMM transitioned into a tail state (via a Poisson jump or an extreme Markov transition), $\theta_{s_t}$ increased, elevating the IV target. This mechanism produced the leverage effect as an emergent property of the return dynamics rather than as a separate model component: large negative returns drove the HMM into low-numbered tail states with $p_{\text{neg}} \approx 0.52$, which in turn elevated $\theta$.

The function $\psi$ captured the shape of the IV surface with five parameters:\label{sec:psi}
\begin{equation}\label{eq:psi}
    \psi(\tau,\; m) = \exp\!\Big(\beta_1 \ln\tau + \beta_2 \ln m + \beta_3 \ln\tau\,\ln m + \beta_4 (\ln m)^2 + \beta_5 (\ln\tau)^2\Big)
\end{equation}
where $\tau = \max(\mathrm{DTE}, 1)$ and $m = K/S$. The parameter $\beta_1$ controlled the linear term-structure decay; $\beta_2$ controlled the skew (negative values produced put skew, with OTM puts having higher IV than OTM calls); $\beta_3$ captured the DTE-skew interaction (the skew flattened at longer maturities); $\beta_4$ controlled the smile curvature; and $\beta_5$ controlled the DTE curvature, capturing the U-shaped ATM term structure where short-dated and long-dated IV were both elevated relative to intermediate maturities.

Rather than specifying a separate initial variance $v_0$, we initialized the variance process at its mean-reversion target:\label{sec:initialization}
\begin{equation}\label{eq:v0}
    v_0 = \theta(s_0,\; \mathrm{DTE},\; K/S_0,\; M_0)
\end{equation}
This ensured that each (ticker, strike, DTE) triple started at its own equilibrium implied volatility. The initial IV surface inherited the smile, skew, and term structure encoded in $\psi$ without requiring any external IV data. Different contracts on the same underlying received different $v_0$ values, and the variance process immediately began mean-reverting around its local $\theta(t)$. The value of this design choice lay in eliminating $v_0$ as a free parameter: the model did not require a separate calibration of the initial IV surface, as it was determined by construction from $\theta$.

The market mood $M_t$ was defined as the fraction of tickers currently occupying tail HMM states, where $N_{\text{tail}}$ denoted the number of extreme states at each end of the state space (inherited from the JumpHMM fit):\label{sec:mood}
\begin{equation}\label{eq:mood}
    M_t = \frac{1}{N_{\text{tickers}}} \sum_{i=1}^{N_{\text{tickers}}} \mathbf{1}\!\left\{s_t^{(i)} \leq N_{\text{tail}} \;\text{or}\; s_t^{(i)} > N - N_{\text{tail}}\right\}
\end{equation}
This signal was fully endogenous to the JumpHMM copula model; it required no external data (e.g., VIX) and introduced no circularity. When multiple tickers simultaneously occupied tail states, the mood rose, elevating $\theta$ and thus IV across all contracts. This mechanism produced correlated IV spikes during market stress events by construction.

The variance process was discretized via Euler-Maruyama with a reflecting boundary at zero and a daily time step $\Delta t = 1/252$:
\begin{equation}\label{eq:euler}
    v_{t+1} = \left| v_t + \kappa\bigl(\theta_t - v_t\bigr)\Delta t + \sigma_v \sqrt{\max(v_t, 0)}\,\sqrt{\Delta t}\, Z_t \right|, \quad Z_t \sim \mathcal{N}(0, 1)
\end{equation}
The reflecting boundary allowed $\theta$ to vary freely across HMM states without requiring the Feller condition $2\kappa\theta > \sigma_v^2$ to hold at every timestep~\cite{andersen2008}. This was a practical necessity, as tail states could drive $\theta$ to values where the Feller condition would otherwise be violated. American option prices were computed from the resulting IV using a CRR binomial tree with 200 steps per contract, which provided sufficient convergence for the DTE range considered.

\subsection{Neural surrogate for $\psi$}\label{sec:psi_nn}

The parametric form of $\psi$ in equation~\eqref{eq:psi} used five shape parameters shared across all contracts on a given ticker. When the model was applied to a broad universe spanning multiple sectors, this global shape became a bottleneck: the semiconductor smile geometry differed systematically from mega-cap ETFs, healthcare names carried idiosyncratic event risk, and energy tickers inherited commodity-driven skew patterns not expressible by a five-term log-polynomial. We therefore replaced the parametric $\psi$ with a neural surrogate~\cite{horvath2021}, keeping the same two inputs $(\ln\tau,\,\ln m)$ and the same multiplicative decomposition $\sigma_{\text{model}}^2 = \theta_{\text{ticker}} \cdot \psi(\ln\tau,\ln m)$, but allowing the shape to be data-driven:
\begin{equation}\label{eq:psi_nn}
    \psi_{\text{NN}}(\ln\tau,\ln m) = \exp\!\bigl(\mathrm{NN}(\tilde z_\tau,\,\tilde z_m)\bigr), \qquad (\tilde z_\tau,\,\tilde z_m) = \bigl((\ln\tau - \mu_\tau)/\sigma_\tau,\; (\ln m - \mu_m)/\sigma_m\bigr)
\end{equation}
where $(\mu_\tau,\sigma_\tau,\mu_m,\sigma_m)$ were global standardization constants computed over the full training set. The network produced $\ln\psi$ so that $\psi > 0$ by construction, preserving the log-Gaussian interpretation of the smile surface used in the parametric form. Each network was a two-hidden-layer multilayer perceptron with $\tanh$ activations: $2 \to 16 \to 16 \to 1$ (337 parameters) for groups with $\geq 2000$ observations and $2 \to 8 \to 8 \to 1$ (105 parameters) for smaller groups, a width chosen to keep the observations-to-parameters ratio between 10 and 50 across all sector fits.

The decomposition defined a natural hierarchy of sharing. At one extreme, a single global $\psi_{\text{NN}}$ shared across all tickers assumed that smile geometry was universal and left only the variance level $\theta_{\text{ticker}}$ to distinguish names. At the other extreme, a per-ticker $\psi_{\text{NN}}$ recovered ticker-specific shape but required enough data per name to resolve the surface without overfitting. We adopted a sector-specific intermediate: one $\psi_{\text{NN}}^{(c)}$ per sector $c$, with all tickers within a sector sharing shape but each ticker retaining its own $\theta_{\text{ticker}}$. This choice reflected the empirical observation that within-sector tickers shared dominant risk factors (technology names tracked semiconductor cycles, healthcare names tracked FDA and trial catalysts, financials tracked rate expectations), so pooling shape across a sector provided a favorable bias-variance tradeoff at the current data volume. As per-ticker sample sizes grew, the same architecture supported relaxation to per-ticker networks without changing the downstream pipeline.

Training minimized the mean squared IV error jointly over the network weights and the per-ticker log-variance levels $\{\ln\theta_{\text{ticker}}\}$:
\begin{equation}\label{eq:nn_loss}
    \min_{w,\,\{\ln\theta_t\}} \;\frac{1}{n} \sum_{i=1}^{n} \Bigl(\sqrt{\theta_{t(i)} \cdot \psi_{\text{NN}}(\ln\tau_i,\ln m_i;\,w)} - \text{IV}_{\text{market},i}\Bigr)^2
\end{equation}
where $t(i)$ denoted the ticker of the $i$th observation. The log-variance levels were initialized from each ticker's mean squared market IV and then updated jointly with the network weights under a single Adam optimizer, which coupled the level and shape fits and avoided the chicken-and-egg problem of fitting one conditional on the other. The learning rate followed a step schedule (1\mbox{e-3} $\to$ 5\mbox{e-4} at epoch 500 $\to$ 2\mbox{e-4} at 1000 $\to$ 1\mbox{e-4} at 1500), and training ran up to 2000 epochs with a best-checkpoint patience of 200 epochs on the training loss. Because the 31-ticker two-day corpus was modest, we did not hold out a temporal validation split at this stage; the reported RMSEs are therefore in-sample fits, and external temporal validation using additional capture days is deferred to future work (Section~\ref{sec:discussion}).
